\documentclass[a4paper,12pt]{article}\usepackage[]{graphicx}\usepackage[]{color}
% maxwidth is the original width if it is less than linewidth
% otherwise use linewidth (to make sure the graphics do not exceed the margin)
\makeatletter
\def\maxwidth{ %
  \ifdim\Gin@nat@width>\linewidth
    \linewidth
  \else
    \Gin@nat@width
  \fi
}
\makeatother

\definecolor{fgcolor}{rgb}{0.345, 0.345, 0.345}
\newcommand{\hlnum}[1]{\textcolor[rgb]{0.686,0.059,0.569}{#1}}%
\newcommand{\hlstr}[1]{\textcolor[rgb]{0.192,0.494,0.8}{#1}}%
\newcommand{\hlcom}[1]{\textcolor[rgb]{0.678,0.584,0.686}{\textit{#1}}}%
\newcommand{\hlopt}[1]{\textcolor[rgb]{0,0,0}{#1}}%
\newcommand{\hlstd}[1]{\textcolor[rgb]{0.345,0.345,0.345}{#1}}%
\newcommand{\hlkwa}[1]{\textcolor[rgb]{0.161,0.373,0.58}{\textbf{#1}}}%
\newcommand{\hlkwb}[1]{\textcolor[rgb]{0.69,0.353,0.396}{#1}}%
\newcommand{\hlkwc}[1]{\textcolor[rgb]{0.333,0.667,0.333}{#1}}%
\newcommand{\hlkwd}[1]{\textcolor[rgb]{0.737,0.353,0.396}{\textbf{#1}}}%
\let\hlipl\hlkwb

\usepackage{framed}
\makeatletter
\newenvironment{kframe}{%
 \def\at@end@of@kframe{}%
 \ifinner\ifhmode%
  \def\at@end@of@kframe{\end{minipage}}%
  \begin{minipage}{\columnwidth}%
 \fi\fi%
 \def\FrameCommand##1{\hskip\@totalleftmargin \hskip-\fboxsep
 \colorbox{shadecolor}{##1}\hskip-\fboxsep
     % There is no \\@totalrightmargin, so:
     \hskip-\linewidth \hskip-\@totalleftmargin \hskip\columnwidth}%
 \MakeFramed {\advance\hsize-\width
   \@totalleftmargin\z@ \linewidth\hsize
   \@setminipage}}%
 {\par\unskip\endMakeFramed%
 \at@end@of@kframe}
\makeatother

\definecolor{shadecolor}{rgb}{.97, .97, .97}
\definecolor{messagecolor}{rgb}{0, 0, 0}
\definecolor{warningcolor}{rgb}{1, 0, 1}
\definecolor{errorcolor}{rgb}{1, 0, 0}
\newenvironment{knitrout}{}{} % an empty environment to be redefined in TeX

\usepackage{alltt}
\usepackage[applemac]{inputenc}
\usepackage[OT1,T1]{fontenc}

\setlength{\textwidth}{160mm}
\setlength{\textheight}{230mm}
\setlength{\oddsidemargin}{-5mm}
\setlength{\topmargin}{-10mm}
% \usepackage{times}
\usepackage{setspace}
\usepackage{lscape}
\usepackage{amsmath,amssymb}
\usepackage{fancyvrb}
\usepackage{multirow}
\usepackage{float}
\usepackage{hyperref}
\usepackage{color}
\usepackage{listings}
\usepackage{fancyhdr}

\usepackage{fancyvrb,moreverb,boxedminipage}


\textwidth=6in
\textheight=8in
\topmargin=0.0in
\oddsidemargin=0in 
\baselineskip=18pt
\headheight=2cm
\setcounter{MaxMatrixCols}{10}

\newcommand{\ns}{\!\!\!}
\newcommand{\e}{\varepsilon}
\newcommand{\be}{\beta}
\newcommand{\s}{\sigma}
\newcommand{\vv}{\vspace{.5cm}}
\DeclareMathOperator{\plim}{plim}
\DefineVerbatimEnvironment{code}{Verbatim}{fontsize=\small}
\DefineVerbatimEnvironment{example}{Verbatim}{fontsize=\small}

\def\by{\mathbf{y}}
\def\bx{\mathbf{x}}
\def\ba{\mathbf{a}}
\def\bb{\mathbf{b}}
\def\be{\mathbf{e}}
\def\bu{\mathbf{u}}
\def\bv{\mathbf{v}}
\def\bz{\mathbf{z}}

%===========
%Greek letter vectors
%===========
\def\bbeta{\mbox{\boldmath $\beta$}}
\def\bgamma{\mbox{\boldmath $\gamma$}}
\def\bvareps{\mbox{\boldmath $\varepsilon$}}
\def\bleta{\mbox{\boldmath $\eta$}}
\def\bmu{\mbox{\boldmath $\mu$}}
\def\btheta{\mbox{\boldmath $\theta$}}
\def\bsigma{\mbox{\boldmath $\sigma$}}
\def\blambgam{\mbox{\boldmath $\lambda_{\gamma}$}}
\def\blambbet{\mbox{\boldmath $\lambda_{\beta}$}}
\def\blambda{\mbox{\boldmath $\lambda$}}

%===========
%Matrices
%===========
\def\bX{\mathbf{X}}
\def\bD{\mathbf{D}(\bsigma)}
\def\bV{\mathbf{V}(\bsigma)}
\def\bVinv{\mathbf{V}^{-1}(\bsigma)}
\def\bR{\mathbf{R}(\bsigma)}
\def\bA{\mathbf{A}}
\def\bB{\mathbf{B}}
\def\bZ{\mathbf{Z}}
\def\bIn{\mathbf{I_{N}}}
\def\bIm{\mathbf{I_{m}}}
\def\bIni{\mathbf{I_{n}}}
\def\SSW{\mathrm{SSW}}
\def\SSB{\mathrm{SSB}}
\def\Normal{\mathcal{N}\!}

%===========
%Parentheses, etc.
%===========
\def\op{\left(}
\def\cp{\right)}
\def\oc{\left[}
\def\cc{\right]}
\def\ok{\left\{}
\def\ck{\right\}}

%===========
%Statistical functions and real numbers.
%===========
\def\exE{\mathbb{E}}
\def\Real{\mathbb{R}}
\def\Var{\mathbb{V}ar}
\def\Proba{\mathbb{P}}

\hypersetup{colorlinks=true,linkcolor=blue,citecolor=blue, urlcolor=blue}
\urlstyle{tt}

\newenvironment{exercices}{\newcounter{moncompteur}\begin{list}
   {\bfseries Ex. \arabic{moncompteur}}{\usecounter{moncompteur}
     \setlength{\labelwidth}{2.6em}\setlength{\leftmargin}{3.1em}}}{\end{list}}
\renewcommand{\labelenumi}{\bfseries\alph{enumi})}
\newcommand{\splus}[3][0]{\texttt{>~#2}\ \ \ \ \textit{#3}\\[#1mm]}

\newenvironment{exo}
{\begin{center}\setlength{\textwidth}{140mm} \underline{Exercise:} \begin{minipage}[t]{120mm}}
{\end{minipage}\setlength{\textwidth}{160mm}\end{center}}


\usepackage{verbatim}

\lstset{ %
  language=R, 
  basicstyle=\ttfamily,
  backgroundcolor=\color{lightgray},   % choose the background color; you must add \usepackage{color} or \usepackage{xcolor}
  xleftmargin= 10 pt,
  xrightmargin= 10 pt,
  escapeinside={\%*}{*)},          % if you want to add LaTeX within your code
  frame=none,                    % adds a frame around the code
  numbers=none,                    % where to put the line-numbers; possible values are (none, left, right)
  numbersep=5pt,                   % how far the line-numbers are from the code
  numberstyle=\tiny\color{mygray}, % the style that is used for the line-numbers
  showspaces=false,                % show spaces everywhere adding particular underscores; it overrides 'showstringspaces'
}

\pagestyle{fancy}

\definecolor{mygreen}{rgb}{0,0.6,0}
\definecolor{mygray}{rgb}{0.5,0.5,0.5}
\definecolor{mymauve}{rgb}{0.58,0,0.82}
\definecolor{lightgray}{gray}{0.95}


%%%%%%%%%%%%%%%%%%%%%%%%%%%%%%%%%%%%%%%%%%%%%%%%%%%%%%%%%%%%%%%%%%%%%%%%%%%%
%

\IfFileExists{upquote.sty}{\usepackage{upquote}}{}
\begin{document}
%

\lhead{University of Geneva\\ \textbf{GLAM}\\
Pr. Eva Cantoni}
\rhead{GSEM\\ Spring 2022 \\ \textbf{Homework 03}}

\begin{center}
\bigskip

\bigskip

\bigskip

\bigskip

\bigskip

\bigskip

\fbox{{\Large Creating functions in \texttt{R}}}

\bigskip

\bigskip

\bigskip
\end{center}

%%%%%%%%%%%%%%%%%%%%%%%%%%%%%%%%%%%%%%%%%%%%%%%%%%%%%%%%%%%%%%%%%%%%%%%%%%%%


\setlength{\parskip}{1ex plus 0.5ex minus 0.2ex}
\parindent=0cm
\setlength{\parskip}{1ex plus 0.5ex minus 0.2ex}
\linespread{1.1}
\sloppy

\section{\underline{Exercise 1}}

In this Exercise we would like you to write your own functions in \texttt{R}. The general procedure for defining new functions is as follows
\begin{center}
\begin{lstlisting}
myfun <- function(argument1,argument2,...){
    expression1
    expression2
}
\end{lstlisting}
\end{center}
where \texttt{myfun} is the name of the function you are creating, \texttt{argument1,argument2,...} can be any type of \texttt{R} objects you are using as input and \texttt{expression} can be any computation and/or already existing function applied to the arguments. By default the result of the last expression is the returned value of the function. Consult the Help manual to learn more about writing your own functions.

\medskip
Write a function (name it \verb"mysummary") to compute and display the average and variance of a (vector) sample $\bx$, using functions such as
\verb"mean()", \verb"var()", \verb"return()" and \verb"list()".

\section{\underline{Exercise 2}}

The goal of this exercise is to find numerical solutions of a set of estimating (score) equations and to compare them with those obtained with the \texttt{R} built-in functions.

Let's consider the \verb"breast-cancer-wisconsin.txt" dataset. We will use the response \verb"status" (coded 0/1) and only one explanatory variable  \verb"nuclei" (plus an intercept). 

We consider the logit function as the link between the random part $\exE[Y_i | \bx_i]$ and the systematic part (linear predictor) $\eta_i=\beta_0 + \beta_1 x_{i1}$.

Based on this, the formula for Newton-Raphson algorithm is:

\begin{align*}
\beta^{(k)}=
\beta^{(k-1)} +& 
\ok 
\sum_{i=1}^n 
\frac{\exp \op - \bx_i^T \beta^{(k-1)} \cp}{\oc 1 + \exp\op -\bx_i\beta^{(k-1)}\cp \cc ^2} 
\bx_i\bx_i^T
\ck ^ {-1}\\
&\times
\ok
\sum_{i=1}^n 
\oc 
y_i - \frac{1}{1 + \exp \op -\bx_i^T\beta^{(k-1)} \cp} 
\cc 
\bx_i 
\ck.
\end{align*}

\begin{enumerate}
  \item{Write an \texttt{R} function that computes an update of the estimated parameters, given the previous value, according to the Newton-Raphson algorithm.}
  \item\label{manually}{Taking the null vector as the starting value for $\bbeta$, perform a few iterations of the Newton-Raphson procedure using a \texttt{while()} in the following way:
  
  \begin{lstlisting}
  while (max(abs(beta0-beta1))>tol & it<maxit){
  ...
  }
\end{lstlisting}
  


where \texttt{beta0} and \texttt{beta1} are the actual and next beta estimations, \texttt{tol} is equal to $10^{-8}$ and \texttt{maxit} is $50$. Report the estimated $\bbeta$ in a table.}
  \item{Write an \texttt{R} function that computes the log-likelihood of your model.}
  \item{Use the built-in function \verb"optim()" to solve the maximization problem (with the ``BFGS'' method). Compare your results with those in \textbf{b)}.}
\end{enumerate}

\section{\underline{Exercise 3}}

%As seen in class, inference on the parameters of a GLM can be performed via the construction of Confidence Intervals by two methods. On one hand, the distribution of $\bbeta$ can be approximated using asymptotic theory on M-estimators. On the other hand, so-called \emph{Profile Likelihood} Confidence Intervals can be created. In \texttt{R}, this procedure is implemented in the \texttt{confint} method.

%The purpose of the current exercise is to understand the construction of Profile Likelihood Confidence Intervals. 

Consider the profile likelihood on page 111 of the course slides\footnote{Based on a handout by Professor David Patterson, University of Montana \url{http://www.math.umt.edu/patterson/}}: 
%
\begin{equation*}
    L_1\op\theta\cp = \underset{\mathbf{\delta}}{\mathrm{max}} \ L\op\theta, \mathbf{\delta} \cp, 
\end{equation*}
%
where $\theta$ is the parameter of interest of dimension 1, and $\delta$ is the
vector of the other model parameters. Under this formulation, we construct the following Likelihood
Ratio (LR) statistic to test $H_0 : \theta = \theta_0$: 
%
\begin{equation*}
    \mathrm{LR}\op \theta_0\cp = 
    2\oc \log L\op \widehat{\theta} , \widehat{\mathbf{\delta}}\cp - \log L\op \theta_0 , \widehat{\mathbf{\delta}}\cp\cc = 
    2 \oc \log L_1\op\widehat{\theta}\cp - \log L_1\op\theta_0\cp \cc \underset{H_0}{\sim}\chi^2_1,
\end{equation*}
%
where $\widehat{\theta}, \widehat{\mathbf{\delta}}$ are the maximum likelihood
(ML) estimators of $\theta$ and $\mathbf{\delta}$. The non-rejection region of
this test is determined by the condition $\mathrm{LR}\op\theta_0\cp < F^{-1}\op
1-\alpha\cp$, where $F$ is the cdf of the $\chi^2_1$ distribution, or
equivalently, all the $\theta_0$ such that: 
%
\begin{equation} \label{eq:nonrejection}
    \log L_1\op \theta_0\cp > \log L\op\widehat{\theta},\widehat{\mathbf{\delta}}\cp - \frac{1}{2}F^{-1}\op 1-\alpha\cp
\end{equation}
%
The points within the \emph{Profile Likelihood} CI will be those $\theta_0$ within the non-rejection region of the LR test.

Consider the goose permit data (file \texttt{goose.txt}) from Bishop and Heberlein (\textit{Amer. J. Agr. Econ.} 61, 1979), where 237 hunters were each offered one of 11 cash amount (bids) ranging from $1\$ $ to $ 200 \$ $ in return of their permits. We use a logit model with \texttt{log(bid)} as the explanatory variable (in addition to the intercept):
$$\log \op \frac{\pi_i}{1-\pi_i} \cp = \beta_0  + \beta_1 \log(\texttt{bid}_i) $$
where $\pi_i$ is the probability that a permit holder would accept bid amount $\texttt{bid}_i$.

%
\begin{enumerate}
    \item Estimate the parameters of the model, extract the log-likelihood with  \texttt{logLik()} to determine the cutoff value as in \eqref{eq:nonrejection} for $\alpha=0.05$ and $\alpha=0.01$. Compute the profiled likelihood CI for $\beta_1$ with the \texttt{confint()} function.
    \item 
\begin{enumerate}
    \item Write an R code to compute the profile likelihood function $L_1(\beta_1)$ (which make use of the estimates of the previous point). 
\item Plot the profile likelihood $L_1(\beta_1)$ and the cutoff values.  
\item Compare the results with the CI computed at point \textbf{a)}.
\end{enumerate}
\end{enumerate}


\end{document}
